% Standalone document. Compile twice with pdflatex for resolved citations.
% Bibliography is embedded; no BibTeX or Biber run is required.
\documentclass[11pt,a4paper]{article}
\usepackage[T1]{fontenc}
\usepackage[utf8]{inputenc}
\usepackage[margin=25mm]{geometry}
\usepackage{lmodern}
\usepackage{microtype}
\usepackage{enumitem}
\usepackage{xurl}
\usepackage[hidelinks]{hyperref}
\setlist[itemize]{leftmargin=*,itemsep=5pt,topsep=5pt}
\setlength{\parindent}{0pt}
\setlength{\parskip}{6pt}
\newcommand{\coeq}{\ensuremath{\mathrm{CO}_{2}\mathrm{e}}}
\newcommand{\entry}[3]{\item \textbf{#1 --- #2.} #3}
\newcommand{\source}[4]{\bibitem{#1} #2. \emph{#3}. \url{#4}. Accessed 24 September 2026.}

\title{Greenhouse Gas and Climate Policy Timeline\\2016--2025}
\author{TRACE3 --- Work Package 4}
\date{}

\begin{document}
\maketitle

This timeline summarizes selected major greenhouse-gas policies, carbon-pricing measures, sustainability directives, and government incentives in the European Union, the United States, and selected Asian economies. Asian coverage includes China, Japan, India, South Korea, Singapore, and Taiwan. US state measures are identified separately. The timeline is selective rather than exhaustive.

\textbf{Classification.} \emph{Regulation} denotes binding legislation, directives, taxes, or mandatory standards, including financial incentives established by law. The term is used broadly, rather than solely in its specific EU legal sense. \emph{Initiative} denotes a government strategy, funding programme, voluntary scheme, or announced commitment. The Paris Agreement is identified separately as an international treaty within the regulation category.

\textbf{Dates and interpretation.} Entries distinguish adoption, announcement, and implementation. Historical descriptions refer to developments during 2016--2025; they do not establish current legal applicability. Some official webpages have subsequently been updated. Carbon taxes and fees, emissions trading, disclosure requirements, emissions standards, and financial incentives should be treated as separate policy mechanisms.

\section*{2016}
\begin{itemize}
\entry{Global}{Regulation / international treaty}{The \textbf{Paris Agreement entered into force on 4 November}, establishing the framework for national climate commitments and reporting. It did not introduce a global carbon tax.\cite{paris}}
\entry{United States}{Regulation / transport emissions}{EPA and NHTSA adopted \textbf{Phase 2 greenhouse-gas and fuel-efficiency standards} for medium- and heavy-duty vehicles.\cite{usphase2}}
\entry{Japan}{Regulation / carbon tax}{The existing climate mitigation tax reached its final scheduled rate in April: \textbf{JPY 289 per tonne of CO\textsubscript{2}}, applied through fossil-fuel taxation.\cite{japantax}}
\end{itemize}

\section*{2017}
\begin{itemize}
\entry{European Union}{Initiative / reporting guidance}{The Commission issued \textbf{non-binding guidance on non-financial reporting}, helping companies disclose environmental and social information under the existing reporting directive.\cite{nfrdguidance}}
\entry{United States / California}{Regulation / carbon trading}{\textbf{AB 398 extended the cap-and-trade framework through 2030}, strengthening long-term carbon-pricing certainty at state level.\cite{california}}
\entry{China}{Initiative / carbon-market development}{The government launched the \textbf{national emissions trading system (ETS) development plan}, initially focused on electricity generation. Actual national trading began in 2021.\cite{china2017,china2025}}
\end{itemize}

\section*{2018}
\begin{itemize}
\entry{European Union}{Regulation / renewable energy and efficiency}{The revised \textbf{Renewable Energy Directive}, energy-efficiency rules, and Energy Union governance framework were adopted. They set original 2030 targets of \textbf{32\% renewable energy and 32.5\% energy-efficiency improvement}.\cite{eu2018}}
\entry{United States}{Regulation / tax incentive}{The \textbf{Bipartisan Budget Act expanded Section 45Q tax credits} for carbon capture, utilization, and storage.\cite{us45q}}
\entry{Singapore}{Regulation / carbon pricing}{The \textbf{Carbon Pricing Act was enacted}, creating the legal basis for a carbon tax starting in 2019.\cite{sgact}}
\end{itemize}

\section*{2019}
\begin{itemize}
\entry{European Union}{Initiative / transition strategy}{The \textbf{European Green Deal launched in December}, setting the direction toward climate neutrality and subsequent energy, industry, transport, and finance reforms.\cite{greendeal}}
\entry{Singapore}{Regulation / carbon-tax implementation}{Carbon taxation began at \textbf{S\$5/t\coeq}, with that introductory rate maintained through 2023.\cite{sgtax}}
\end{itemize}

\section*{2020}
\begin{itemize}
\entry{European Union}{Regulation / sustainable finance}{The \textbf{Taxonomy Regulation entered into force}, establishing a framework for classifying environmentally sustainable economic activities.\cite{taxonomy}}
\entry{United States}{Regulation / industrial greenhouse gases}{The \textbf{American Innovation and Manufacturing Act} mandated an \textbf{85\% phasedown in HFC production and consumption by 2036}, relative to statutory baselines.\cite{aim}}
\entry{China}{Initiative / national climate commitment}{China announced targets to \textbf{peak CO\textsubscript{2} emissions before 2030 and achieve carbon neutrality before 2060}.\cite{chinapledge}}
\entry{Japan}{Initiative / industrial decarbonization}{The \textbf{Green Growth Strategy launched in December}, outlining industrial and technology pathways toward carbon neutrality by 2050.\cite{japangrowth}}
\end{itemize}

\section*{2021}
\begin{itemize}
\entry{European Union}{Regulation / binding climate targets}{The \textbf{European Climate Law} established climate neutrality by \textbf{2050} and at least \textbf{55\% net GHG reduction by 2030}, compared with 1990.\cite{climatelaw}}
\entry{European Union}{Regulation / green investment funding}{The \textbf{Recovery and Resilience Facility} required at least \textbf{37\% of each national recovery plan's allocation} to support climate objectives.\cite{rrf}}
\entry{United States}{Regulation / public investment}{The \textbf{Bipartisan Infrastructure Law} provided more than \textbf{US\$62 billion to the Department of Energy} for clean-energy infrastructure, demonstrations, and related programmes.\cite{bil}}
\entry{China}{Regulation / carbon-market implementation}{\textbf{National ETS trading began in July}, initially covering the power-generation sector.\cite{china2025}}
\entry{Japan}{Initiative / technology funding}{The \textbf{Green Innovation Fund began supporting long-term decarbonization projects}, backed initially by \textbf{JPY 2 trillion} under the FY2020 supplementary budget.\cite{japanfund}}
\end{itemize}

\section*{2022}
\begin{itemize}
\entry{European Union}{Regulation / corporate disclosure}{The \textbf{Corporate Sustainability Reporting Directive (CSRD) was adopted}, expanding and standardizing sustainability reporting, with phased implementation.\cite{csrd}}
\entry{European Union}{Initiative / energy transition}{\textbf{REPowerEU launched in May}, accelerating energy savings, renewable deployment, and diversification away from Russian fossil fuels.\cite{repower}}
\entry{United States}{Regulation / financial incentives}{The \textbf{Inflation Reduction Act (IRA) was enacted}, introducing or expanding tax incentives for clean electricity, vehicles, manufacturing, hydrogen, energy efficiency, and carbon capture.\cite{ira}}
\entry{South Korea}{Regulation / climate governance}{The \textbf{Carbon Neutrality Framework Act and implementing decree became effective in March}, establishing the legal framework for the country's carbon-neutral transition.\cite{korea}}
\entry{India}{Regulation / carbon-market enabling law}{The \textbf{Energy Conservation Amendment Act} empowered the government to establish a carbon-credit trading scheme.\cite{indiaccts}}
\end{itemize}

\section*{2023}
\begin{itemize}
\entry{European Union}{Regulation / carbon pricing}{The \textbf{EU ETS reform was adopted}, strengthening the emissions cap, extending coverage to shipping from 2024, and establishing a separate future system for fuels used in buildings and road transport.\cite{fit55}}
\entry{European Union}{Regulation / carbon border adjustment}{The \textbf{Carbon Border Adjustment Mechanism (CBAM) entered its transitional reporting phase on 1 October} for selected carbon-intensive imports. \textbf{2023--2025 involved reporting}; the financial regime was scheduled from 2026.\cite{cbam}}
\entry{European Union}{Regulation / emissions disclosure}{The \textbf{European Sustainability Reporting Standards (ESRS) were adopted}. ESRS E1 covers Scope 1, 2, and 3 emissions, subject to materiality and applicable phase-ins.\cite{esrs}}
\entry{European Union}{Regulation / renewable energy}{The revised \textbf{Renewable Energy Directive (RED III) raised the binding EU 2030 renewable-energy target to 42.5\%}, with an additional ambition to reach 45\%.\cite{red3}}
\entry{United States}{Regulation / methane}{EPA finalized stronger \textbf{oil and gas methane standards}, including requirements addressing existing sources nationwide.\cite{methane2023}}
\entry{India}{Regulation / carbon trading}{The \textbf{Carbon Credit Trading Scheme was notified}, establishing compliance and offset mechanisms. This established the framework rather than marking a fully operational trading market.\cite{indiaccts}}
\entry{India}{Initiative / production incentives}{The \textbf{National Green Hydrogen Mission was approved}, with an initial \textbf{INR 19,744 crore} allocation, including incentives for hydrogen production and electrolyser manufacturing.\cite{indiah2}}
\entry{Japan}{Regulation / transition finance}{The \textbf{GX Promotion Act was enacted}, establishing the framework for green-transition financing and phased carbon pricing.\cite{gxact}}
\entry{Japan}{Initiative / voluntary carbon trading}{The \textbf{GX League's trial emissions trading system started in April}, with voluntary company participation.\cite{gxleague}}
\end{itemize}

\section*{2024}
\begin{itemize}
\entry{European Union}{Regulation / shipping emissions}{\textbf{Maritime transport entered the EU ETS}, with allowance obligations phased in, starting with \textbf{40\% of covered 2024 emissions}.\cite{fit55}}
\entry{European Union}{Regulation / supply-chain due diligence}{The \textbf{Corporate Sustainability Due Diligence Directive (CSDDD) was adopted}, covering environmental and human-rights impacts and, as adopted, climate-transition planning requirements. Implementation was scheduled for later years.\cite{csddd}}
\entry{European Union}{Regulation / clean manufacturing}{The \textbf{Net-Zero Industry Act was adopted}, supporting clean-technology manufacturing through measures including streamlined permitting and procurement provisions.\cite{nzia}}
\entry{United States}{Regulation / power-sector emissions}{EPA finalized \textbf{CO\textsubscript{2} standards for existing coal-fired and new gas-fired power plants}.\cite{power2024}}
\entry{Singapore}{Regulation / carbon-tax increase}{The carbon tax increased from \textbf{S\$5 to S\$25/t\coeq\ for 2024--2025}.\cite{sgtax}}
\end{itemize}

\section*{2025}
\begin{itemize}
\entry{European Union}{Regulation / reporting delays}{The \textbf{``stop-the-clock'' directive was adopted}, postponing CSRD application for the second and third reporting waves by two years, and the CSDDD transposition deadline and first application phase by one year.\cite{stopclock}}
\entry{European Union}{Initiative / industrial investment}{The \textbf{Clean Industrial Deal launched}, followed in June by a state-aid framework facilitating support for clean energy, industrial decarbonization, and clean-technology manufacturing.\cite{cleandeal}}
\entry{United States}{Regulation / incentive rollback}{The \textbf{One Big Beautiful Bill Act} ended federal new, used, and commercial clean-vehicle credits for vehicles acquired \textbf{after 30 September 2025}.\cite{ev2025}}
\entry{United States}{Regulation / methane-charge rollback}{The implementing \textbf{Waste Emissions Charge rule was disapproved under the Congressional Review Act}, and EPA removed it from the regulations in May.\cite{wec}}
\entry{China}{Regulation / carbon-market expansion}{The national ETS expanded to \textbf{steel, cement, and aluminium smelting}, extending carbon-pricing obligations beyond electricity generation.\cite{china2025}}
\entry{Taiwan}{Regulation / carbon fee}{Carbon fees began accruing on \textbf{2025 emissions}, at a standard \textbf{NT\$300/t\coeq}, with preferential rates for approved reduction plans. First payments were scheduled for May 2026.\cite{taiwan}}
\end{itemize}

\clearpage
\begin{thebibliography}{99}
\small
\source{paris}{UNFCCC}{The Paris Agreement}{https://unfccc.int/process-and-meetings/the-paris-agreement}
\source{usphase2}{US Environmental Protection Agency (2016)}{Final Rule for Phase 2 Greenhouse Gas Emissions Standards and Fuel Efficiency Standards for Medium- and Heavy-Duty Engines and Vehicles}{https://www.epa.gov/regulations-emissions-vehicles-and-engines/final-rule-phase-2-greenhouse-gas-emissions-standards}
\source{japantax}{Ministry of the Environment, Japan}{Tax for Climate Change Mitigation}{https://www.env.go.jp/content/900453366.pdf}
\source{nfrdguidance}{European Commission (2017; updated 2019)}{Commission guidelines on non-financial reporting}{https://finance.ec.europa.eu/publications/commission-guidelines-non-financial-reporting_en}
\source{california}{California Air Resources Board (2018)}{CARB provides further flexibility and cost-containment for companies in cap-and-trade program}{https://ww2.arb.ca.gov/news/carb-provides-further-flexibility-and-cost-containment-companies-cap-and-trade-program}
\source{china2017}{National Development and Reform Commission, China (2017)}{Notice issuing the national carbon emissions trading market construction plan for the power-generation sector [Chinese]}{https://zfxxgk.ndrc.gov.cn/web/iteminfo.jsp?id=2944}
\source{china2025}{Ministry of Ecology and Environment, China (2025)}{Questions and answers on the work plan for national carbon-market coverage of steel, cement, and aluminium smelting [Chinese]}{https://www.mee.gov.cn/ywdt/zbft/202503/t20250326_1104767.shtml}
\source{eu2018}{Council of the European Union (2018)}{Energy efficiency, renewables, governance of the Energy Union: Council signs off on 3 major clean energy files}{https://www.consilium.europa.eu/en/press/press-releases/2018/12/04/energy-efficiency-renewables-governance-of-the-energy-union-council-signs-off-on-3-major-clean-energy-files/pdf/}
\source{us45q}{US Department of Energy (2020)}{DOE Loan Guarantees and IRS Guidance on 45Q Tax Credit Could Benefit Carbon Capture Projects}{https://www.energy.gov/edf/articles/doe-loan-guarantees-irs-guidance-45q-tax-credit-could-benefit-carbon-capture-projects}
\source{sgact}{Singapore Statutes Online}{Carbon Pricing Act 2018}{https://sso.agc.gov.sg/Act/CPA2018}
\source{greendeal}{European Commission (2019)}{Communication on the European Green Deal}{https://commission.europa.eu/publications/communication-european-green-deal_en}
\source{sgtax}{Ministry of Sustainability and the Environment, Singapore}{Carbon Pricing Act}{https://www.mse.gov.sg/policies/climate-change/carbon-pricing-act/}
\source{taxonomy}{European Commission}{EU taxonomy for sustainable activities}{https://finance.ec.europa.eu/sustainable-finance/tools-and-standards/eu-taxonomy-sustainable-activities_en}
\source{aim}{US Environmental Protection Agency}{Frequent Questions on the Phasedown of Hydrofluorocarbons}{https://www.epa.gov/hfcs/frequent-questions-phasedown-hydrofluorocarbons}
\source{chinapledge}{State Council of the People's Republic of China (2021)}{China's new five-year blueprint paves way for 2060 carbon-neutrality}{https://english.www.gov.cn/news/topnews/202103/09/content_WS6046cf92c6d0719374afa6a5.html}
\source{japangrowth}{Ministry of Economy, Trade and Industry, Japan}{Green Growth Strategy Through Achieving Carbon Neutrality in 2050}{https://www.meti.go.jp/english/policy/energy_environment/global_warming/ggs2050/index.html}
\source{climatelaw}{European Commission}{European Climate Law}{https://climate.ec.europa.eu/areas-action/european-climate-law_en}
\source{rrf}{European Commission}{Finance and the Green Deal}{https://commission.europa.eu/strategy-and-policy/priorities-2019-2024/european-green-deal/finance-and-green-deal_en}
\source{bil}{US Department of Energy (2021)}{President Biden Signs Historic Bipartisan Infrastructure Deal to Usher in the Clean Energy Future}{https://www.energy.gov/policy/articles/president-biden-signs-historic-bipartisan-infrastructure-deal-usher-clean-energy}
\source{japanfund}{Ministry of Economy, Trade and Industry, Japan (2021)}{Green Growth Strategy Through Achieving Carbon Neutrality in 2050: Overview}{https://www.meti.go.jp/english/policy/energy_environment/global_warming/ggs2050/pdf/ggs_overview_all.pdf}
\source{csrd}{Council of the European Union (2022)}{Council gives final green light to corporate sustainability reporting directive}{https://www.consilium.europa.eu/en/press/press-releases/2022/11/28/council-gives-final-green-light-to-corporate-sustainability-reporting-directive/}
\source{repower}{European Commission}{REPowerEU}{https://commission.europa.eu/topics/energy/repowereu_en}
\source{ira}{US Internal Revenue Service}{Credits and deductions under the Inflation Reduction Act of 2022}{https://www.irs.gov/credits-and-deductions-under-the-inflation-reduction-act-of-2022}
\source{korea}{Ministry of Environment, Republic of Korea (2022)}{Press release on the entry into force of the Carbon Neutrality Act and enforcement decree, 25 March}{https://www.me.go.kr/eng/web/board/read.do?boardId=1516150&boardMasterId=522&menuId=461}
\source{indiaccts}{Press Information Bureau, Government of India (2024)}{Government statement on the Indian Carbon Market and Carbon Credit Trading Scheme, 9 December}{https://www.pib.gov.in/PressReleasePage.aspx?PRID=2082528&lang=1&reg=3}
\source{fit55}{Council of the European Union (2023)}{Fit for 55: Council adopts key pieces of legislation delivering on 2030 climate targets}{https://www.consilium.europa.eu/en/press/press-releases/2023/04/25/fit-for-55-council-adopts-key-pieces-of-legislation-delivering-on-2030-climate-targets/}
\source{cbam}{European Commission}{Carbon Border Adjustment Mechanism}{https://taxation-customs.ec.europa.eu/carbon-border-adjustment-mechanism_en}
\source{esrs}{European Commission (2023)}{Commission Delegated Regulation (EU) 2023/2772 on sustainability reporting standards}{https://eur-lex.europa.eu/legal-content/en/TXT/?uri=CELEX:32023R2772}
\source{red3}{Council of the European Union (2023)}{Renewable energy: Council adopts new rules}{https://www.consilium.europa.eu/en/press/press-releases/2023/10/09/renewable-energy-council-adopts-new-rules/}
\source{methane2023}{US Environmental Protection Agency (2023)}{EPA's Final Rule to Reduce Methane and Other Harmful Pollution from Oil and Natural Gas Operations and Related Actions}{https://www.epa.gov/controlling-air-pollution-oil-and-natural-gas-operations/epas-final-rule-reduce-methane-and-other}
\source{indiah2}{Press Information Bureau, Government of India (2023)}{Cabinet approves National Green Hydrogen Mission}{https://www.pib.gov.in/PressReleaseIframePage.aspx?PRID=1888547&lang=2&reg=48}
\source{gxact}{Ministry of Economy, Trade and Industry, Japan}{GX Promotion Strategy}{https://www.meti.go.jp/policy/energy_environment/global_warming/transition/pathways_to_green_transformation_eng.pdf}
\source{gxleague}{Agency for Natural Resources and Energy, Japan (2023)}{Japan's energy policy toward achieving GX (Part 2)}{https://www.enecho.meti.go.jp/en/category/special/article/detail_179.html}
\source{csddd}{Council of the European Union (2024)}{Corporate sustainability due diligence: Council gives its final approval}{https://www.consilium.europa.eu/en/press/press-releases/2024/05/24/corporate-sustainability-due-diligence-council-gives-its-final-approval/}
\source{nzia}{Council of the European Union (2024)}{Industrial policy: Council gives final approval to the net-zero industry act}{https://www.consilium.europa.eu/en/press/press-releases/2024/05/27/industrial-policy-council-gives-final-approval-to-the-net-zero-industry-act/}
\source{power2024}{US Environmental Protection Agency (2024)}{Final Carbon Pollution Standards: Overview Fact Sheet}{https://www.epa.gov/system/files/documents/2024-04/cps-111-fact-sheet-overview.pdf}
\source{stopclock}{Council of the European Union (2025)}{Simplification: Council gives final green light on the stop-the-clock mechanism}{https://www.consilium.europa.eu/en/press/press-releases/2025/04/14/simplification-council-gives-final-green-light-on-the-stop-the-clock-mechanism-to-boost-eu-competitiveness-and-provide-legal-certainty-to-businesses/}
\source{cleandeal}{European Commission (2025)}{Clean Industrial Deal State Aid Framework (CISAF)}{https://competition-policy.ec.europa.eu/about/contribution-clean-just-and-competitive-transition/clean-industrial-deal-state-aid-framework-cisaf_en}
\source{ev2025}{US Internal Revenue Service}{Clean vehicle tax credits}{https://www.irs.gov/clean-vehicle-tax-credits}
\source{wec}{US Environmental Protection Agency}{Methane Emissions Reduction Program and GHGRP Subpart W}{https://www.epa.gov/inflation-reduction-act/methane-emissions-reduction-program-and-ghgrp-subpart-w-petroleum-and}
\source{taiwan}{Ministry of Environment, Taiwan (2024)}{Taiwan's Ministry of Environment Announces Fee-Charging Rates of Carbon Fees}{https://www.moenv.gov.tw/en/news/press-releases/31552.html}
\end{thebibliography}
\end{document}
